% Generated draft: requires mathematical and source review.
\documentclass[11pt]{article}
\usepackage[margin=25mm]{geometry}
\usepackage{fontspec}
\setmainfont{Latin Modern Roman}
\usepackage{amsmath,amssymb,mathtools,unicode-math}
\setmathfont{Latin Modern Math}
\usepackage{enumitem,needspace,xurl,hyperref}
\setlength{\parindent}{0pt}
\setlength{\parskip}{5pt}
\setlength{\emergencystretch}{3em}
\newcommand{\N}{\mathbb{N}}
\newcommand{\Z}{\mathbb{Z}}
\newcommand{\Q}{\mathbb{Q}}
\newcommand{\R}{\mathbb{R}}
\newcommand{\C}{\mathbb{C}}
\newcommand{\F}{\mathbb{F}}
\newcommand{\A}{\mathbb{A}}
\newcommand{\PP}{\mathbb{P}}
\newcommand{\E}{\mathbb{E}}
\newcommand{\eps}{\varepsilon}
\newcommand{\dd}{\,\mathrm{d}}
\begin{document}
\section*{UnsolvedMath ID: 30000124}
\textbf{Source problem code:} catalogue entry
\subsection{Definitions and mathematical statement}
Let \(X\) be a smooth complex surface, meaning a complex manifold of dimension two, and let \(E\to X\) be a differentiable \emph{complex} vector bundle of rank \(r\geq1\). A holomorphic structure is an integrable \((0,1)\)-operator
\[
 \overline\partial_E:C^\infty(E)\to\Omega^{0,1}(E),\qquad
 \overline\partial_E(fs)=\overline\partial f\otimes s+f\overline\partial_Es,
 \qquad\overline\partial_E^{\,2}=0.
\]
The source asks which underlying differentiable bundles admit such an operator, not for uniqueness or a moduli classification of all holomorphic structures on one bundle.

Its Chern-data version asks, for each rank, to describe
\[
 \mathcal H_r(X)=\{(c_1(F),c_2(F)):F\text{ a holomorphic rank-}r\text{ vector bundle on }X\}
 \subseteq H^2(X;\Z)\times H^4(X;\Z).
\]
Here \(c_2=0\) for line bundles. On a compact connected surface \(c_2\) may be identified with its Chern number, but in general it is a cohomology class. Any such bundle has \(c_1(F)\) in the image of \(c_1:\operatorname{Pic}(X)\to H^2(X;\Z)\), through its determinant; that necessary condition is not asserted sufficient for all ranks.

The source's ``equivalently'' uses the standard classification of complex bundles in real dimension four by rank and the relevant Chern classes, with the usual rank-one convention. That topological identification should be matched to the intended smooth surface setting. It does not identify different holomorphic bundles having the same Chern data. Compactness, projectivity, K\"ahlerness, and surface class are not stated in the original; changing them changes the existence problem and must not be hidden.

The supplied literature triage dated 21 August 2026 reports a Schwarzenberger criterion in the algebraic-surface case and rank-two existence results for several non-K\"ahler surfaces. It supplies neither the full criterion nor the exact surface/rank hypotheses, and describes the general arbitrary-surface, all-ranks problem as incomplete. No unprovided numerical existence criterion or universal extension of those special cases is asserted here.
\subsection{Short English statement}
Characterize differentiable complex bundles on a complex surface that admit an integrable holomorphic structure, or their realizable rank/Chern data. Compactness and surface class matter. The supplied curation reports algebraic and special rank-two results without a complete all-surface, all-rank criterion.
\subsection{Sources}
\begin{itemize}
\item Local UnsolvedMath registry, permanent ID \(30000124\), exact source code \texttt{catalogue entry}, Oberwolfach Reports Open Problems. Statement status \texttt{corrected\_verified}; verification restores variable notation and removes a redundant label.
\item ``Komplexe Analysis (2005),'' Oberwolfach Reports, pp. 2171--2216, DOI \texttt{10.4171/owr/2004/42}: \url{https://doi.org/10.4171/owr/2004/42}. Proposed year and DOI refer to 2004, while the citation label says 2005.
\item Supplied ``Oberwolfach Report 42, holomorphic bundles on surfaces'': \url{https://ems.press/content/serial-article-files/45960}.
\item Supplied label ``Teleman, holomorphic vector bundles on non-algebraic surfaces'': \url{https://arxiv.org/abs/math/0201236}; supplied ``Teleman, rank-two bundles on non-Kaehler elliptic surfaces'': \url{https://arxiv.org/abs/math/0306191}. Full theorem details and the stated Schwarzenberger criterion are not reproduced.
\item Literature checked 21 August 2026; supplied status \texttt{partially\_solved}. Registry created 25 August 2026 and updated 16 September 2026. This dated curation triage requires independent verification before citation.
\item Registry curation attribution: Ulam AI / UnsolvedMath Contributors, CC BY 4.0. Underlying report and literature retain their supplied provenance; this statement-only draft requires human review.
\end{itemize}
\end{document}
